\documentclass{IEEEtran}

\usepackage{cite}
\usepackage{amsmath,amssymb}
\usepackage{booktabs}
\usepackage[hidelinks]{hyperref}
\usepackage{balance}

\begin{document}

\title{Evolutionary Test Generation with Diversity-Aware Mutation Fitness}

\author{
\IEEEauthorblockN{Minwook Lee, Minwoo Kim, Junho Hwang, and Seonwoo Ji}
\IEEEauthorblockA{CS453 Team 6\\
KAIST}
}

\maketitle

\begin{abstract}
EvoSuite optimises mutation score by counting killed mutants, without considering which mutation operators they belong to.
Two test suites can therefore achieve the same score while covering very different fault models.
We extended EvoSuite's strong-mutation fitness with Shannon entropy over killed-mutant operator types and coupled this diversity signal into the base objective through several modes (additive, multiplicative, exponential, and capped-multiplicative).
A Python automation pipeline ran approximately 3{,}790 EvoSuite experiments across 14 Apache Commons and Guava classes in four phases.
After fixing a population-size issue that blocked evolution in early phases, tuned diversity coupling improved median mutation score on 12 of 13 evaluated classes and jointly improved mutation score and operator entropy on 9 of 13.
Effects are class-specific; \texttt{CharUtils} was the sole regression.
Overall, the results partially support our hypothesis that diversity-aware selection can broaden fault-model exploration, but a single global setting does not work everywhere.
Code: \href{https://github.com/minux-lee/CS453-team6}{github.com/minux-lee/CS453-team6}.
\end{abstract}

\section{Introduction}

Automated unit test generation is a practical way to improve software quality.
EvoSuite~\cite{fraser2011evosuite,fraser2014mutation} uses an evolutionary algorithm to generate JUnit tests that maximise coverage and mutation-based criteria.
The strong-mutation objective minimises branch distance plus per-mutant infection/propagation distances~\cite{fraser2014mutation}.
Once a mutant is killed, the operator type does not affect fitness further.

This creates a blind spot: two suites with the same mutation score can kill very different mixes of operator types.
For example, one suite might kill mostly arithmetic mutants while another spreads kills across arithmetic, comparison, and condition-negation mutants.
The second suite covers a broader fault model but receives identical selection pressure.

Our hypothesis is that coupling operator-type diversity into fitness steers search toward broader exploration and improves mutation score without sacrificing it.
This report describes the problem, our EvoSuite extension, experiments, and a before/after comparison against a baseline with diversity coupling disabled.

\section{Problem and Hypothesis}

Mutation testing measures how many artificial faults a suite detects~\cite{papadakis2019mutation}.
EvoSuite's strong-mutation fitness combines branch coverage with distances to alive mutants (lower is better).

The score counts \emph{how many} mutants are killed, not \emph{which} operator types they represent.
EvoSuite uses nine primary operators (field deletion, statement deletion, unary insertion, variable replacement, constant replacement, arithmetic replacement, comparison replacement, bitwise replacement, condition negation) plus an ``other'' bucket.

\textbf{Hypothesis:} a normalised Shannon-entropy measure over killed-mutant operator types, when coupled into fitness, improves mutation score and operator diversity together compared with the stock objective under the same time and population budget.

We use operator entropy as a proxy for fault-model breadth and do not validate against real faults (e.g., Defects4J).

\section{Approach}

\subsection{Diversity Metric}

Let $c_i$ be killed mutants in operator bucket $i$, $p_i = c_i / \sum_j c_j$, and $n$ the number of buckets present:
\begin{align}
v &= -\textstyle\sum_{i=1}^{n} p_i \log p_i, \label{eq:entropy-a}\\
\hat{v} &= v / \log n \in [0,1], \qquad \delta = 1 - \hat{v}. \label{eq:entropy-b}
\end{align}
$\hat{v}$ is normalised entropy; $\delta$ is \emph{diversity deficiency} (0 = diverse, 1 = single-type concentration).
Entropy is computed over killed mutants on each fitness evaluation.

\subsection{Coupled Fitness}

EvoSuite minimises fitness. The base value is
$F_{\mathrm{base}} = f_{\mathrm{branch}} + \sum_{m \in \mathrm{alive}} d_{\min}(m)$.
Table~\ref{tab:fitness} lists coupling modes controlled by \texttt{-Ddiversity\_coupling}, \texttt{-Ddiversity\_k}, and \texttt{-Ddiversity\_cap}.

\begin{table}[htbp]
\caption{Coupled fitness modes ($\delta = 1-\hat{v}$).}
\label{tab:fitness}
\centering
\small
\begin{tabular}{@{}ll@{}}
\toprule
Mode & Formula \\
\midrule
\texttt{NONE} & $F_{\mathrm{base}}$ \\
\texttt{MULT} & $F_{\mathrm{base}}(1 + k\delta)$ \\
\texttt{ADD} & $F_{\mathrm{base}} + k\delta$ \\
\texttt{EXP} & $F_{\mathrm{base}}\, e^{k\delta}$ \\
\texttt{CAP\_MULT} & $F_{\mathrm{base}}(1 + \min(k\delta, c))$ \\
\bottomrule
\end{tabular}
\end{table}

Plain multiplicative coupling at large $k$ can collapse mutation score by over-penalising low diversity; \texttt{CAPPED\_MULTIPLICATIVE} limits the multiplier with cap~$c$.

\section{Implementation}

We started from the latest EvoSuite release on GitHub, integrated our changes into the project repository, and built with Maven.
Table~\ref{tab:files} lists the main modified files.

\begin{table}[htbp]
\caption{EvoSuite source changes.}
\label{tab:files}
\centering
\footnotesize
\begin{tabular}{@{}p{0.48\columnwidth}p{0.46\columnwidth}@{}}
\toprule
File & Change \\
\midrule
\texttt{StrongMutation\-SuiteFitness.java} &
Operator bucketing, entropy, coupled fitness, stats export \\
\texttt{Properties.java} &
\texttt{DiversityCoupling} enum; \texttt{diversity\_k}, \texttt{diversity\_cap} \\
\texttt{RuntimeVariable.java} &
\texttt{MutantTypeEntropy}, \texttt{MutantTypeEntropyNorm} \\
\bottomrule
\end{tabular}
\end{table}

\textbf{\texttt{StrongMutationSuiteFitness.java}.}
We restructured the evaluator to separate base mutation distances from diversity coupling, map mutants to operator buckets, and export entropy to \texttt{statistics.csv}.
Seven defects in the initial skeleton (null-pointer risks, NaN propagation, inverted optimisation direction) were fixed during integration.

The repository also contains a Python \texttt{automation/} package (benchmark registry, parallel EvoSuite runner, resumable sweeps, phase drivers) and an \texttt{analysis/} package (CSV parsing, delta metrics, plotting).

\section{Experimental Setup}

\subsection{Benchmarks}

We used 14 utility classes from commons-math3, commons-lang3, commons-text, commons-codec, and Guava.
Table~\ref{tab:benchmarks} lists the benchmarks and domains.

\begin{table}[htbp]
\caption{Benchmark classes (14 total).}
\label{tab:benchmarks}
\centering
\footnotesize
\setlength{\tabcolsep}{3pt}
\begin{tabular}{@{}lll@{}}
\toprule
Class & Library & Dom. \\
\midrule
ArithmeticUtils, Fraction, & & \\
Precision, Complex & math3 & num. \\
IntMath & Guava & num. \\
IEEE754rUtils, NumberUtils, & & \\
CharUtils, BooleanUtils & lang3 & mixed \\
WordUtils, LevenshteinDist. & text & str. \\
Strings & Guava & str. \\
Hex, Soundex & codec & enc. \\
\bottomrule
\end{tabular}
\end{table}

\subsection{Phases and Budgets}

Experiments proceeded in four phases (Table~\ref{tab:phases}).
Phase~A used a coarse $k$ grid $\{0.25, 0.5, 1.0, 2.0, 4.0, 8.0\}$, population 50, 40\,s budget, and 3 seeds.
Phase~B re-ran promising configurations with 10 seeds.
Phases~C and~D are the primary evidence: population 5, denser $k$ grids, and 8 or 6 seeds.

\begin{table}[htbp]
\caption{Experimental phases.}
\label{tab:phases}
\centering
\small
\begin{tabular}{@{}cccc@{}}
\toprule
Ph. & \#Cls & Pop. & Budget \\
\midrule
A & 14 & 50 & 40\,s \\
B & 14 & 50 & 40\,s \\
C & 8 & 5 & 60\,s \\
D & 6 & 5 & 120\,s \\
\bottomrule
\end{tabular}
\end{table}

\textbf{Population artefact.}
With population 50 and short budgets, most Phase~A/B runs reported \texttt{Generations=0}: the initial random population exhausted the time limit before evolution occurred.
Reducing population to 5 in Phases~C/D enabled real search ($\sim$35 generations on average); all comparative claims below refer to Phase~C/D.

\textbf{Control settings.}
All runs used \texttt{STRONGMUTATION}, \texttt{-Dtest\_archive=false}, and \texttt{NONE} as baseline.
Metrics: mutation score (\texttt{Coverage}), normalised entropy (\texttt{MutantTypeEntropyNorm}), and deltas vs.\ matched baseline.

\section{Results}

\subsection{Phase A/B Screening}

Phase~A/B established directional trends:
\texttt{EXPONENTIAL} coupling failed to train reliably and was excluded from later phases.
\texttt{ADDITIVE} and \texttt{MULTIPLICATIVE} with $k \in [0, 1.5]$ showed stable gains; plain \texttt{MULTIPLICATIVE} with $k > 2$ frequently collapsed mutation score, motivating \texttt{CAPPED\_MULTIPLICATIVE}.

\subsection{Phase C (Smaller Classes)}

Phase~C evaluated eight classes (\texttt{Soundex}, \texttt{Hex}, \texttt{Fraction}, \texttt{Precision}, \texttt{NumberUtils}, \texttt{IEEE754rUtils}, \texttt{CharUtils}, \texttt{Strings}) with a dense $k$ grid, eight seeds, and population 5.
Table~\ref{tab:phc_k} lists the best $k$ per coupling mode selected by average mutation score.

\begin{table}[htbp]
\caption{Phase C: best $k$ per coupling mode.}
\label{tab:phc_k}
\centering
\small
\begin{tabular}{@{}lc@{}}
\toprule
Coupling & Best $k$ \\
\midrule
\texttt{ADDITIVE} & 0.5 \\
\texttt{MULTIPLICATIVE} & 0.2 \\
\texttt{CAPPED\_MULTIPLICATIVE} & 8.0 (cap 1.0) \\
\bottomrule
\end{tabular}
\end{table}

\texttt{CharUtils} was an outlier: the baseline already achieved high mutation score and entropy; diversity pressure shifted selection toward operator balance at the cost of fewer killed mutants ($\Delta\mathrm{MS} \approx -0.054$).

\subsection{Phase D (Larger Classes)}

Phase~D targeted six larger classes (\texttt{ArithmeticUtils}, \texttt{Complex}, \texttt{BooleanUtils}, \texttt{WordUtils}, \texttt{Fraction}, \texttt{IntMath}) with a 120\,s budget.
Table~\ref{tab:phd_k} lists the best $k$ values.

\begin{table}[htbp]
\caption{Phase D: best $k$ per coupling mode.}
\label{tab:phd_k}
\centering
\small
\begin{tabular}{@{}lc@{}}
\toprule
Coupling & Best $k$ \\
\midrule
\texttt{ADDITIVE} & 0.5 \\
\texttt{MULTIPLICATIVE} & 0.75 \\
\texttt{CAPPED\_MULTIPLICATIVE} & 2.0 (cap 1.0) \\
\bottomrule
\end{tabular}
\end{table}

\subsection{Combined Phase C+D Summary}

Across 13 unique classes covered by Phases~C and~D (8 + 6 minus overlapping \texttt{Fraction}):
\textbf{12/13} improved median mutation score under the best tuned configuration per class;
\textbf{9/13} improved both mutation score and normalised entropy at that configuration;
\textbf{1/13} (\texttt{CharUtils}) regressed.

Table~\ref{tab:gains} highlights notable per-class gains at best tuned settings.

\begin{table}[htbp]
\caption{Selected per-class improvements (Phase C+D).}
\label{tab:gains}
\centering
\small
\setlength{\tabcolsep}{3pt}
\begin{tabular}{@{}llcr@{}}
\toprule
Class & Coupling & $k$ & $\Delta$MS \\
\midrule
Fraction & MULT & 0.75 & +.091 \\
Precision & ADD & 1.5 & +.079 \\
Soundex & MULT & 0.5 & +.044 \\
ArithmeticUtils & CAP\_MULT & 0.5 & +.051 \\
CharUtils & --- & --- & $-.054$ \\
\bottomrule
\end{tabular}
\end{table}

Domain-level aggregation (Table~\ref{tab:domain}) shows mutation score improvements in numeric, boolean, and encoding domains; the string domain regressed only on \texttt{CharUtils}.

\begin{table}[htbp]
\caption{Domain-level trend at best configuration (Phase C+D).}
\label{tab:domain}
\centering
\small
\begin{tabular}{@{}ll@{}}
\toprule
Domain & Trend \\
\midrule
Numeric & Improved on most classes \\
Boolean & Improved \\
Encoding & Improved (\texttt{Soundex}, \texttt{Hex}) \\
String & Improved except \texttt{CharUtils} \\
\bottomrule
\end{tabular}
\end{table}

With $\leq 8$ seeds per setting, Mann--Whitney $U$ tests ($p < 0.05$) rarely reached significance; we emphasise median effect sizes.

\subsection{Before/After Interpretation}

Relative to the \texttt{NONE} baseline under identical budgets:

\begin{enumerate}
\item Modest $k$ (0.2--0.75) most often \emph{increased} mutation score rather than trading it for entropy alone.
\item \texttt{MULTIPLICATIVE} with $k \geq 2$ reliably \emph{degraded} mutation score.
\item \texttt{CAPPED\_MULTIPLICATIVE} stabilised high-$k$ behaviour (e.g., \texttt{ArithmeticUtils}).
\item Effects are \emph{conditional}: no single $(\mathrm{mode}, k)$ dominated all classes.
\end{enumerate}

Table~\ref{tab:recommend} summarises practical configuration choices from our experiments.

\begin{table}[htbp]
\caption{Practical configuration recommendations.}
\label{tab:recommend}
\centering
\footnotesize
\begin{tabular}{@{}lll@{}}
\toprule
Use case & Mode & $k$ \\
\midrule
Safe default & \texttt{ADDITIVE} & 0.5--1.0 \\
Strongest gains & \texttt{MULTIPLICATIVE} & 0.5--0.75 \\
High $k$ without collapse & \texttt{CAP\_MULT} & 2--4 (cap 1.0) \\
Avoid & \texttt{MULTIPLICATIVE} & $\geq 2$ \\
\bottomrule
\end{tabular}
\end{table}

\section{Limitations and Future Work}

Phases~C/D used population~5 and 60--120\,s budgets; results may differ with EvoSuite defaults and longer runs.
\texttt{test\_archive=false} isolates suite-fitness effects but lowers absolute scores.
Operator entropy is a proxy, not validated against real faults~\cite{just2014mutants}.
Future work: larger budgets, per-generation tracking, phased objectives, and Defects4J evaluation.

\section{Conclusion}

We extended EvoSuite strong-mutation fitness with operator-type entropy coupling.
On 13 benchmark classes, tuned settings improved mutation score on 12/13 and jointly improved score and entropy on 9/13.
The hypothesis is partially supported: diversity-aware fitness helps on many classes but needs per-class tuning.
A practical default is \texttt{ADDITIVE} with $k \approx 0.5$--1.0; use \texttt{CAPPED\_MULTIPLICATIVE} when stronger diversity pressure is needed.

\balance
\begin{thebibliography}{00}

\bibitem{fraser2011evosuite}
G.~Fraser and A.~Arcuri, ``EvoSuite: Automatic test suite generation for Java classes,'' in \emph{Proc. ISSTA}, 2011.

\bibitem{fraser2014mutation}
G.~Fraser and A.~Arcuri, ``Achieving higher test quality with mutation-based test generation,'' \emph{IEEE Trans. Softw. Eng.}, vol.~40, no.~9, pp.~1041--1059, 2014.

\bibitem{papadakis2019mutation}
M.~Papadakis \emph{et al.}, ``Mutation testing advances,'' \emph{Adv. Comput.}, vol.~112, pp.~1--75, 2019.

\bibitem{just2014mutants}
R.~Just \emph{et al.}, ``Are mutants a valid substitute for real faults?'' in \emph{Proc. FSE}, 2014.

\end{thebibliography}

\end{document}
